% !TEX root = ../main.tex

\chapter{Validazione preliminare dello strumento}
\label{cap:risultati}

Questo capitolo riporta la validazione dello strumento presentato nel
Capitolo~\ref{cap:metodo} e stabilisce in che misura
l'obiettivo della tesi sia stato raggiunto. Le evidenze sono di tre tipi e
vanno lette in modo diverso. Le coppie di casi di prova dicono se ogni modulo
riconosce la condizione per cui è stato costruito. Il codice non scritto per i
moduli, cioè i repository esterni e crAPI, mostra come si comporta lo strumento
quando la condizione non è stata preparata per lui. Un insieme di prova di tre applicazioni scritte da terzi,
infine, è stato riservato alla validazione e analizzato con tutti i moduli in
un'unica campagna. Dopo aver richiamato obiettivi e impostazione della
validazione, il capitolo presenta i risultati dei due tipi di prova, li discute
per categoria alla luce della classificazione, risponde alle domande di ricerca
e si chiude con i limiti e le minacce alla validità.

\section{Obiettivi e impostazione della validazione}
\label{sec:ris-oggetto}

La validazione deve fornire le evidenze sperimentali per rispondere alle due
domande di ricerca: determinare quale combinazione di livelli di analisi consenta
la copertura delle categorie tecniche (\RQ{1.1}) e chiarire, per le categorie
ibride, quale quota rimanga riconoscibile staticamente e quale sia la natura del
residuo (\RQ{1.2}).

Come riassunto nella Tabella~\ref{tab:ris-evidenze}, le evidenze raccolte si
articolano su tre livelli complementari:
\begin{enumerate}
  \item \textbf{Casi di prova dedicati:} per ciascun modulo dello strumento è stata
    predisposta una coppia di scenari a specchio (una versione vulnerabile e una
    irrobustita), coprendo sia l'infrastruttura IaC sia i servizi API sviluppati
    nei framework supportati (Flask per Python ed Express per Node.js). Questa verifica
    controlla puntualmente la capacità di ciascuna regola di rilevare la debolezza
    target senza generare falsi allarmi;
  \item \textbf{Codice esterno non dedicato:} repository pubblici ed esempi realistici
    preesistenti (incluso il benchmark OWASP crAPI~\cite{owasp2026crapi}) consentono
    di osservare il comportamento dello strumento su basi di codice non realizzate ad hoc;
  \item \textbf{Insieme di prova indipendente:} una campagna di valutazione complessiva
    condotta su tre applicazioni terze indipendenti --- sadcloud~\cite{nccgroup2021sadcloud}
    per l'infrastruttura Terraform, VAmPI~\cite{erev0s2026vampi} e DVAPI~\cite{payatu2025dvapi}
    per le API. In questa campagna globale tutti i moduli operano congiuntamente, e
    l'accuratezza è misurata in termini di \emph{recall} (generale e nel perimetro)
    e \emph{precisione minima}.
\end{enumerate}

\begin{table}[htbp]
  \centering
  \small
  \linespread{1.0}\selectfont
  \begin{tabularx}{\linewidth}{@{}>{\raggedright\arraybackslash}p{0.27\linewidth}>{\raggedright\arraybackslash}X>{\raggedright\arraybackslash}X>{\raggedright\arraybackslash}p{0.14\linewidth}@{}}
    \toprule
    \textbf{Modulo} & \textbf{Casi di prova} & \textbf{Codice non scritto per il modulo} & \textbf{Insieme di prova} \\
    \midrule
    IaC misconfiguration & manifest Terraform a protezione crescente & crAPI & sadcloud \\
    API1 -- BOLA & coppia in Flask ed Express & repository etichettato; crAPI & VAmPI, DVAPI \\
    API2 -- Broken Authentication & coppia gemella & -- & VAmPI, DVAPI \\
    API3 -- BOPLA & coppia sul profilo utente & crAPI & VAmPI, DVAPI \\
    API4 -- Unrestricted Resource Consumption & coppia gemella & -- & VAmPI, DVAPI \\
    API5 -- BFLA & coppia con le sette debolezze & repository con specifica OpenAPI & VAmPI, DVAPI \\
    API7 -- SSRF & coppia gemella & -- & VAmPI, DVAPI \\
    API8 -- Security Misconfiguration & coppia gemella & -- & VAmPI, DVAPI \\
    API10 -- Unsafe Consumption & coppia gemella & bersaglio realistico & VAmPI, DVAPI \\
    \bottomrule
  \end{tabularx}
  \caption[Prove di cui si riportano gli esiti]{Prove di cui si riportano gli
    esiti, per modulo. Un trattino indica che per quel modulo non si riportano
  esiti su codice non scritto per il modulo.}
  \label{tab:ris-evidenze}
\end{table}

Per i casi di prova il riferimento di verità è noto per costruzione; per crAPI
è stato ricostruito a mano su una revisione fissata, individuando nel codice i
punti in cui si manifestano le vulnerabilità documentate dal progetto. Una
segnalazione conta come corrispondente a una voce del riferimento quando
riguarda la stessa categoria OWASP e lo stesso punto di esposizione
dell'interfaccia, identificato da metodo e percorso normalizzato, qualunque sia
il livello di analisi che l'ha prodotta; questa convenzione è stata fissata
prima di osservare i risultati.

A completamento dell'infrastruttura di test, una suite di prove d'unità automatizzate
garantisce la robustezza e la correttezza logica dei singoli moduli e delle euristiche
di analisi prima dell'esecuzione sui bersagli reali.

\section{Risultati sui casi di prova}
\label{sec:ris-casi}

L'analisi dei casi di prova dedicati ha lo scopo di verificare la correttezza
funzionale dei singoli moduli in condizioni controllate, confermando sia la
capacità di rilevare le violazioni introdotte, sia l'assenza di falsi allarmi
sulle corrispondenti versioni sicure.

Sul fronte infrastrutturale, l'analisi dei manifest Terraform (Tabella~\ref{tab:iac-ranking})
conferma l'accuratezza del modello di rischio: le configurazioni aperte o prive di
autorizzazione si collocano sistematicamente ai vertici della graduatoria di severità,
mentre le risorse protette ottengono punteggi minimi legati a soli irrobustimenti.

L'ordinamento dipende dal catalogo delle policy, che assegna natura e severità
a ciascun controllo fallito di Checkov secondo le regole della
Tabella~\ref{tab:catalogo-policy}, provate in ordine: la mappatura esplicita,
che nella configurazione usata copre 99 controlli, poi il prefisso e le parole
chiave, e infine un ripiego predefinito che lascia la segnalazione non
classificata con severità \texttt{MEDIUM}. La regola applicata determina anche
la confidenza del Finding, massima per la mappatura esplicita e via via minore
per le successive.

\begin{table}[htbp]
  \centering
  \small
  \linespread{1.0}\selectfont
  \begin{tabularx}{\linewidth}{@{}c l >{\hsize=0.95\hsize}X >{\hsize=1.05\hsize}X@{}}
    \toprule
    \textbf{\#} & \textbf{Regola} & \textbf{Quando si applica} &
    \textbf{Esito assegnato (esempio)} \\
    \midrule
    1 & Mappatura esplicita & l'identificativo del controllo è elencato nel
    catalogo & natura, severità e flag presi dal catalogo
    \par\smallskip \texttt{CKV\_AWS\_20} $\rightarrow$ esposizione,
    \texttt{CRITICAL}, pubblica \\
    \addlinespace
    2 & Prefisso & l'identificativo inizia con una famiglia nota &
    esito della famiglia
    \par\smallskip \texttt{CKV\_SECRET\_*} $\rightarrow$ esposizione,
    \texttt{HIGH}, dati sensibili \\
    \addlinespace
    3 & Parole chiave & il nome ufficiale del controllo contiene termini noti
    (i gruppi difensivi hanno la precedenza su quelli di esposizione) &
    esito del gruppo
    \par\smallskip \emph{logging}, \emph{versioning} $\rightarrow$
    irrobustimento, \texttt{LOW} \\
    \addlinespace
    4 & Predefinito & nessuna delle regole precedenti si applica &
    non classificato, \texttt{MEDIUM} \\
    \bottomrule
  \end{tabularx}
  \caption[Risoluzione dei controlli nel catalogo delle policy]{Le quattro
    regole vengono provate nell'ordine della colonna~\#: la prima che si applica
    assegna natura, severità e contesto al controllo fallito di Checkov, e
    determina anche la confidenza. Le esposizioni hanno severità predefinita
    \texttt{HIGH}, gli irrobustimenti \texttt{LOW}, salvo indicazione esplicita
  nel catalogo.}
  \label{tab:catalogo-policy}
\end{table}

\begin{table}[htbp]
  \centering
  \begin{tabular}{@{}lcl@{}}
    \toprule
    \textbf{Risorsa} & \textbf{Punteggio} & \textbf{Natura} \\
    \midrule
    ACL pubblica e policy aperta & 9,2 & esposizione \\
    Chiave cablata (Lambda) & 8,8 & esposizione \\
    Policy IAM \texttt{*} su \texttt{*} & 8,6 & esposizione \\
    Metodi senza autorizzazione & 7,4 & esposizione \\
    Blocchi di accesso disattivati & 6,8 & esposizione \\
    Bucket sicuri (due) & 3,2 & solo irrobustimenti \\
    \bottomrule
  \end{tabular}
  \caption[Punteggi di rischio sul bersaglio Terraform]{Punteggi di rischio
  delle voci del bersaglio Terraform, in ordine decrescente.}
  \label{tab:iac-ranking}
\end{table}

Per i moduli applicativi, la Tabella~\ref{tab:ris-gemelle} riassume gli esiti ottenuti
sulle coppie gemelle. I moduli statici (BOLA, BFLA, BOPLA, SSRF, Resource Consumption,
Security Misconfiguration e Unsafe Consumption) rilevano la totalità delle debolezze
introdotte senza produrre alcun falso positivo sulle versioni irrobustite. Nel caso della
BOPLA si osserva inoltre la cooperazione tra i livelli: le segnalazioni dedotte dal contratto
vengono confermate o ridimensionate in base al comportamento osservato a runtime. Per la
Broken Authentication, le prove dinamiche identificano puntualmente le violazioni attive
senza fallire sulle implementazioni corrette.

\begin{table}[htbp]
  \centering
  \small
  \linespread{1.0}\selectfont
  \begin{tabularx}{\linewidth}{@{}l>{\hsize=0.9\hsize}X>{\hsize=1.1\hsize}Xcc@{}}
    \toprule
    \textbf{Modulo} & \textbf{Versione vulnerabile} & \textbf{Versione corretta} & \textbf{VP} & \textbf{FP} \\
    \midrule
    BOLA (livello statico) & 12 casi su 12 segnalati & nessuna segnalazione & 12 & 0 \\
    Broken Authentication & 5 prove su 22 falliscono & nessuna prova fallisce & -- & -- \\
    BOPLA & S01, S02 e S04 segnalate e confermate a runtime & resta S04, vera sul contratto e smentita a runtime & 3 & 0 \\
    Resource Consumption & 9 casi su 9 segnalati & nessuna segnalazione & 9 & 0 \\
    BFLA & 7 debolezze su 7 segnalate & nessuna segnalazione & 7 & 0 \\
    SSRF & 3 casi su 3 segnalati & nessuna segnalazione & 3 & 0 \\
    Security Misconfiguration & 7 casi su 7 segnalati & nessuna segnalazione & 7 & 0 \\
    Unsafe Consumption & 3 regole su 3 segnalate & nessuna segnalazione & 3 & 0 \\
    \bottomrule
  \end{tabularx}
  \caption[Esiti sulle coppie gemelle]{Esiti sulle coppie gemelle. VP e FP
    indicano i veri e i falsi positivi; per la Broken Authentication, le cui
    prove sono dinamiche, si riporta il numero di prove fallite. Il bersaglio
    infrastrutturale, valutato
    sull'ordinamento dei punteggi, è riportato nella
  Tabella~\ref{tab:iac-ranking}.}
  \label{tab:ris-gemelle}
\end{table}

Questi esiti confermano che ciascun modulo risponde ai requisiti teorici in ambiente ideale;
la reale efficacia dello strumento e i suoi confini applicativi emergono dal confronto con
codice non scritto ad hoc, discusso nelle sezioni seguenti.

\subsection{Risultati su codice non scritto per i moduli}

\subsubsection{Repository esterni}

Sul repository esterno, l'analisi statica della BOLA ha prodotto una segnalazione OA-001 per ciascuno dei quattro metodi della risorsa vulnerabile, senza segnalazioni sulla risorsa protetta. Uno dei metodi vulnerabili non era inoltre presente nella specifica OpenAPI, mostrando come l'analisi statica possa individuare casi che non sarebbero stati raggiunti dalla sola analisi dinamica.

È stata rilevata anche un'ulteriore situazione potenzialmente riconducibile a una BOLA, relativa a un endpoint di importazione che permette a un utente autenticato di modificare il proprietario di un progetto. Il caso non è presente tra le vulnerabilità dichiarate dal repository e non è stato verificato durante la prova dinamica, quindi non viene considerato né un vero né un falso positivo.

L'analisi dinamica della BOLA ha eseguito 108 prove e ha confermato gli accessi non autorizzati alla risorsa vulnerabile tramite GET, PUT e PATCH. Non sono invece emerse segnalazioni sulla risorsa protetta. La cancellazione della risorsa vulnerabile, individuata dall'analisi statica, non è stata confermata.

Nel complesso, l'analisi statica ha individuato quattro metodi vulnerabili, mentre quella dinamica ne ha confermati tre.

Sul repository accompagnato dalla specifica OpenAPI, l'analisi della BFLA ha prodotto tre segnalazioni BF-006 relative a endpoint privi di autenticazione e controllo del ruolo. Gli endpoint sono però utilizzati dal progetto per preparare e ripristinare i dati delle prove automatiche e non rappresentano quindi vulnerabilità dell'applicazione. Gli endpoint destinati al normale utilizzo dell'applicazione non hanno prodotto segnalazioni.

La prova ha inoltre evidenziato un problema nella visualizzazione del metodo HTTP: per alcune segnalazioni viene indicato GET anche quando l'endpoint utilizza un metodo diverso. Il percorso rimane corretto, ma il metodo mostrato può risultare errato.

\subsubsection{crAPI}

Su crAPI, i risultati variano in base alla categoria analizzata. Per la parte infrastrutturale, le 240 segnalazioni iniziali sono state ridotte a 46 dopo il filtraggio e il raggruppamento. Sono stati individuati sia problemi reali sia un falso positivo relativo a credenziali che non appartengono a crAPI.

Per la BOLA, l'analisi statica non ha prodotto segnalazioni. La BOPLA ha invece riconosciuto alcune proprietà protette, ma non tutti i casi presenti nell'applicazione.

Il confronto con le vulnerabilità documentate da crAPI, riportato nella Tabella~\ref{tab:ris-crapi}, mostra che nessuna delle tredici sfide riconducibili alla Top 10 API è stata rilevata. Il risultato evidenzia quindi i limiti della rilevazione statica quando viene applicata a un'applicazione reale con strutture e modalità di implementazione diverse da quelle considerate nello sviluppo delle regole.

\begin{table}[htbp]
  \centering
  \small
  \linespread{1.0}\selectfont
  \begin{tabularx}{\linewidth}{@{}>{\raggedright\arraybackslash}p{0.32\linewidth}l>{\raggedright\arraybackslash}p{0.17\linewidth}X@{}}
    \toprule
    \textbf{Sfida di crAPI} & \textbf{Cat.} & \textbf{Esito} & \textbf{Motivo} \\
    \midrule
    1 -- veicolo di un altro utente & API1 & non segnalata & servizio Java, fuori dal livello statico; livello dinamico non eseguito su crAPI \\
    2 -- report dei meccanici altrui & API1 & non segnalata & vista Django definita come classe, identificativo nella query \\
    4, 5 -- dati altrui e proprietà interne esposti & API3 & non segnalate & nessuna segnalazione statica \\
    8--10 -- mass assignment su ordini, saldo e video & API3 & non segnalate & campi assegnati uno per uno, fuori dalla regola S01 \\
    3, 14, 15 -- reset della password, accesso senza autenticazione, JWT falsificabili & API2 & non valutate & esiti del modulo non riportati \\
    6 -- DoS tramite ``contact mechanic'' & API4 & non valutata & esiti del modulo non riportati \\
    7 -- cancellazione di un video altrui & API5 & non valutata & modulo non eseguito su crAPI \\
    11 -- richiesta verso un host esterno & API7 & non valutata & esiti del modulo non riportati \\
    12, 13, 16--18 -- iniezioni e chatbot & -- & fuori perimetro & non appartengono alla Top 10 API \\
    \bottomrule
  \end{tabularx}
  \caption[Riferimento di verità di crAPI ed esiti dello strumento]{Sfide
    documentate da crAPI alla revisione~1.1.5, categoria OWASP ed esito dello
  strumento.}
  \label{tab:ris-crapi}
\end{table}

Sul bersaglio realistico le regole dell'Unsafe Consumption non producono
segnalazioni, un esito che riflette le caratteristiche architetturali del
codice esaminato più che un difetto di copertura. UC-002 e UC-003 identificano
il servizio esterno a partire da un indirizzo statico presente nel sorgente,
mentre lì gli endpoint delle terze parti sono definiti solo a runtime e restano
fuori dal perimetro dell'analisi statica. UC-001, che non dipende
dall'indirizzo ma segue il flusso del dato dalla risposta della terza parte
fino al punto di riuso, non trova punti di iniezione non protetti, come query
costruite per concatenazione o oggetti creati da dizionari non filtrati: manca
il bersaglio dell'iniezione differita. Lo stesso codice resta coperto, da
un'angolazione complementare, dal modulo per la Server-Side Request Forgery
(Sottosezione~\ref{sec:rilevamento-ssrf}). Proprio perché il difetto non c'è, però,
questo bersaglio non dice nulla sulla sensibilità del modulo, che resta
documentata dalle sole applicazioni gemelle.

\subsection{Insieme di prova}
\label{sec:ris-prova}

La Tabella~\ref{tab:ris-prova} riporta, per ciascun bersaglio, le voci della
scansione statica e le segnalazioni di ogni modulo. VAmPI compare due volte
perché può essere avviata in modalità vulnerabile o sicura con lo stesso
codice.

\begin{table}[htbp]
  \centering
  \small
  \linespread{1.0}\selectfont
  \begin{tabular}{@{}lcccc@{}}
    \toprule
    \textbf{Modulo} & \textbf{VAmPI vuln.} & \textbf{VAmPI sicuro} & \textbf{DVAPI} & \textbf{sadcloud} \\
    \midrule
    Scansione statica (voci) & 28 & 28 & 73 & 38 \\
    BOLA, livello dinamico & 1 su 12 prove & 0 su 12 prove & non eseguita & -- \\
    BOLA, livello statico & 0 & 0 & 0 & 0 \\
    Broken Authentication & 1 & 1 & non eseguita & -- \\
    BOPLA & 0 & 0 & 0 & 0 \\
    Resource Consumption & 3 & 3 & 3 & 0 \\
    BFLA & 0 & 0 & 1 & 0 \\
    SSRF & 10 & 10 & 0 & 0 \\
    Security Misconfiguration & 5 & 5 & 3 & 1 \\
    Unsafe Consumption & 0 & 0 & 0 & 0 \\
    \bottomrule
  \end{tabular}
  \caption[Esiti sull'insieme di prova]{Voci della scansione statica e
    segnalazioni per modulo sui bersagli dell'insieme di prova. Un trattino
    indica che il modulo non si applica: sadcloud non ha un'applicazione in
  esecuzione.}
  \label{tab:ris-prova}
\end{table}

Il risultato più chiaro viene da VAmPI: l'analisi statica restituisce le
stesse segnalazioni nelle due modalità, e solo il livello dinamico della BOLA
le distingue. Nella modalità vulnerabile conferma che un utente può leggere il
libro di un altro (\texttt{GET /books/v1/\{titolo\}}), la vulnerabilità
documentata da VAmPI; nella modalità sicura le stesse richieste vengono
respinte. La differenza sta nel comportamento, non nel codice, e nessuna
analisi statica potrebbe vederla. Su DVAPI la BOLA dinamica e la Broken
Authentication non sono state eseguite per due limiti noti dello strumento: il
formato dei token di DVAPI e la posizione del suo file delle dipendenze.

La Tabella~\ref{tab:ris-metriche} confronta le segnalazioni con le
vulnerabilità documentate da ciascun progetto. La \emph{recall} è la quota di
vulnerabilità trovate; la \emph{recall nel perimetro} esclude quelle che lo
strumento non può trovare per costruzione, ad esempio perché nessun modulo ne
copre la categoria. La \emph{precisione minima} è la quota di segnalazioni che
corrispondono a una vulnerabilità documentata: è una stima per difetto, perché
conta come errate anche le segnalazioni non ancora esaminate a mano.

\begin{table}[htbp]
  \centering
  \small
  \linespread{1.0}\selectfont
  \begin{tabular}{@{}lccccccc@{}}
    \toprule
    \textbf{Bersaglio} & \textbf{Voci} & \textbf{VP} & \textbf{FN} & \textbf{Fuori perim.} & \textbf{Recall} & \textbf{Recall nel perim.} & \textbf{Prec. minima} \\
    \midrule
    VAmPI & 8 & 2 & 3 & 3 & 25\% & 40\% & 11\% \\
    DVAPI & 9 & 1 & 5 & 3 & 11\% & 17\% & 2\% \\
    sadcloud & 84 & 51 & 33 & 0 & 61\% & 61\% & 25\% \\
    \bottomrule
  \end{tabular}
  \caption[Metriche di accuratezza sull'insieme di prova]{Metriche di
    accuratezza sull'insieme di prova. Le voci fuori dalla Top 10 (una per VAmPI,
  una per DVAPI) sono escluse.}
  \label{tab:ris-metriche}
\end{table}

I risultati sono molto diversi fra infrastruttura e API. Su sadcloud Checkov
trova 51 delle 84 misconfiguration documentate: le trova quando il problema
riguarda una singola risorsa, come un database non cifrato o un bucket
pubblico, e le perde quando riguarda l'account nel suo insieme, come l'assenza
di CloudTrail. La precisione minima si ferma al 25\% perché Checkov segnala
anche molti irrobustimenti che il riferimento non elenca.

Sulle API lo strumento trova poco. Delle tre corrispondenze, solo una
riconosce davvero la vulnerabilità: la BOLA di VAmPI confermata dal livello
dinamico. Le altre due coincidono con la voce documentata per categoria ed
endpoint, ma Spectral vi segnala un problema diverso. Restano mancati, fra gli
altri, il mass assignment di VAmPI e, su DVAPI, il caricamento di file senza
limiti e la SSRF. Le precisioni minime sulle API sono molto basse perché gran
parte delle segnalazioni viene dalle regole di Spectral sul contratto OpenAPI
e, non essendo ancora stata esaminata a mano, conta come errata.

\section{Discussione per categoria}
\label{sec:ris-categorie}

\subsection{Categorie ibride}
\label{sec:ris-ibride}

Le prove confermano la divisione del lavoro prevista dalla classificazione. La
parte statica trova la condizione necessaria e raggiunge punti che la verifica
dinamica non vedrebbe: i quattro veri positivi della BOLA su operazioni che il
contratto dichiara protette, dove Spectral tace, e il metodo vulnerabile del
repository esterno assente dalla specifica. La parte dinamica, dove c'è,
restituisce ciò che il codice non può dire, come mostra la segnalazione S04
della BOPLA, vera sul contratto e smentita dal sistema in esecuzione. Il
livello dinamico della BOLA lo mostra su due bersagli in esecuzione: sul
repository etichettato distingue la risorsa vulnerabile da quella protetta, e su
VAmPI distingue due modalità che hanno lo stesso codice, cioè proprio ciò che
l'analisi statica non può dire. Il meccanismo con cui conferme e smentite
raggiungono la stessa voce del motore di correlazione resta invece verificato
soltanto dalle prove d'unità, e la segnalazione sull'endpoint di importazione,
che la scansione dinamica non ha raggiunto, è l'esempio concreto di un residuo
dinamico rimasto aperto.

Nella BFLA il residuo resta aperto per costruzione, e le prove mostrano quanto
sarebbe difficile chiuderlo. OWASP giudica facile la rilevabilità della
categoria, perché le API sono strutturate e il modo di raggiungere una funzione
è prevedibile~\cite{owasp2023api5}, ma il giudizio riguarda metà del problema:
formulare la chiamata candidata è facile, stabilire se la risposta fosse dovuta
non lo è, perché un \texttt{200 OK} attesta che l'operazione è stata eseguita,
non che fosse legittimo eseguirla. Il contributo della letteratura non consiste
nell'ispezionare il codice, ma nel ricostruire la specifica mancante,
ricavandola dalle pagine raggiungibili da ciascun ruolo~\cite{sun2011static},
dalla coerenza con cui i controlli sono applicati dentro una partizione per
ruolo~\cite{son2011rolecast} o dal contesto di
autorizzazione~\cite{monshizadeh2014mace}; un'alternativa rinuncia a inferire e
scrive regole di sicurezza per le API REST, verificandole con fuzzing
stateful~\cite{atlidakis2020checking}, mentre un lavoro recente affronta lo
stesso problema in modalità black box~\cite{liu2025bacscan}. Le tre
segnalazioni BF-006 sul repository esterno illustrano il punto: la regola
stabilisce che quelle funzioni sono raggiungibili senza ruolo, non se debbano
esserlo, e a deciderlo è chi conosce lo scopo degli endpoint.

Nella Broken Authentication il residuo si vede nella forma stessa degli esiti.
Sulle applicazioni gemelle le prove distinguono la versione vulnerabile da
quella irrobustita, ma circa metà resta inconcludente: una prova dinamica può
accertare un difetto solo dove il meccanismo esiste e una sessione è
disponibile. Su DVAPI il modulo non ha raggiunto la parte dinamica, e i suoi
limiti sono discussi nella Sezione~\ref{sec:ris-limiti}.

\subsection{Categorie tecniche}
\label{sec:ris-statiche}

Per l'Unrestricted Resource Consumption, classificata come tecnica, le
applicazioni gemelle danno un esito pieno, ma il limite principale del modulo è
di natura metodologica e merita una discussione a parte. Il modulo non vede l'API in
funzione: quanto costi davvero un endpoint, o se un tetto dichiarato venga poi
rispettato, si scopre solo mettendo il sistema sotto carico e osservandone i
tempi di risposta e le chiamate verso l'esterno. Le due fonti primarie sembrano
qui contraddirsi, perché OWASP giudica facile la rilevabilità, dato che
richieste con parametri di dimensione crescente e l'osservazione di stato, tempo
e lunghezza della risposta bastano a far emergere il
problema~\cite{owasp2023api4}, mentre CWE-770 giudica di utilità limitata
l'analisi statica per le allocazioni prive di limite~\cite{cwe770}. La
contraddizione è apparente: OWASP descrive un tester in \emph{black
box}, CWE uno strumento che legge il codice. Sul versante dinamico
l'evidenza è consolidata, dalla generazione di molte richieste in un intervallo
breve al rilevamento per iniezione di
carico~\cite{cwe770,antunes2008resource}; il modulo, lavorando per indizi, sa
dire che un limite manca, ma non se la sua assenza basti a mettere in difficoltà
il servizio. Una sua segnalazione va letta quindi come una protezione da
aggiungere, e la severità dice quanto quella protezione conti, non quanto sia
probabile un attacco.

L'Unsafe Consumption presenta una copertura mista. L'analisi statica ne
individua la \emph{predisposizione} --- la validazione assente, il canale in
chiaro, la redirezione seguita ciecamente --- mentre la verifica della
sfruttabilità effettiva richiede un piano dinamico che metta alla prova il
sistema in esecuzione contro un fornitore ostile e controllato: un banco che
osservi il traffico in uscita e risponda di proposito con dati malevoli o con
una redirezione verso un host pericoloso. Una verifica del genere è possibile
solo se il servizio reale consente di immettervi contenuti, oppure sostituendolo
con un doppio in collaudo, che però verifica l'applicazione e non
l'integrazione reale. Il residuo non coperto è quindi di natura dinamica, nel
senso di \RQ{1.2}, e la sua verifica andrebbe integrata con quella della
Server-Side Request Forgery. Per la Server-Side Request Forgery e la Security
Misconfiguration gli esiti pieni sulle coppie gemelle confermano che la
condizione si legge negli artefatti. Sull'insieme di prova, però, nessuna delle
due trova la voce documentata di DVAPI: né la richiesta verso un URL arbitrario
né la traccia dello stack restituita con un token non valido. I loro limiti
sono discussi nella Sezione~\ref{sec:ris-limiti}.

\subsection{Categorie senza modulo}
\label{sec:ris-scoperte}

Due categorie restano prive di un modulo, e per ragioni diverse. Per API6:2023
il rilevamento automatico si scontra con l'assenza di un
oracolo: nessun file analizzabile, che sia codice, contratto o
configurazione, risponde alla domanda se un dato flusso in una data azienda
vada protetto, e la mancanza di una mappatura CWE registra il fatto che non
esiste una debolezza del software da nominare. Le tecniche note per i difetti
di logica applicativa cercano deviazioni da un comportamento
atteso~\cite{felmetsger2010logic,pellegrino2014logic}, mentre qui il flusso
viene percorso correttamente e l'attacco sta nel ripeterlo. Restano le
contromisure comportamentali a runtime indicate da OWASP~\cite{owasp2023api6},
che riconoscono l'automazione durante l'attacco anziché la vulnerabilità prima
che avvenga. Nei termini del criterio, la categoria ricade quindi fra quelle
di logica di business.

API9:2023 resta invece fuori dalla tripartizione, perché presuppone che il
sistema da valutare non sia noto (Sottosezione~\ref{sec:criterio}). OWASP ne
giudica media la rilevabilità senza indicare un metodo, limitandosi a osservare
che una documentazione obsoleta rende più difficile individuare e correggere le
vulnerabilità~\cite{owasp2023api9}: il difetto ostacola il proprio rilevamento,
perché la documentazione che servirebbe a trovarlo è proprio ciò che manca. Per via
automatica si può soltanto confrontare rappresentazioni diverse dello stesso
sistema: le rotte implementate nel codice ma non dichiarate nel contratto sono
candidate a essere API ombra, quelle dichiarate ma non implementate segnalano
documentazione obsoleta, e nessuna delle due è una vulnerabilità. Segnali di
questo tipo compaiono nel sistema solo come effetto collaterale di altri moduli,
nell'handler della BOLA assente dal contratto e nella funzione ombra della BFLA,
e non costituiscono una copertura della categoria. Il confronto ha del resto un
limite intrinseco, perché opera fra due rappresentazioni della stessa
applicazione e non dice nulla sugli host su cui essa è distribuita: il primo
scenario OWASP è precisamente il caso in cui codice e contratto coincidono e il
difetto sta nella topologia di rete, il che richiede ricognizione
infrastrutturale e non analisi del software.

\section{Risposta alle domande di ricerca}
\label{sec:ris-rq}

\paragraph{Risposta a \RQ{1.1}}
Per le categorie tecniche la combinazione di livelli non è un raffinamento, ma
la condizione perché un giudizio sia possibile. Sulla parte infrastrutturale il
manifest da solo non basta: servono il catalogo delle policy, che assegna a ogni
controllo natura e severità, e la correlazione, che riunisce i controlli sulla
stessa risorsa. Così le 86 segnalazioni del bersaglio Terraform si ordinano come
atteso, e su crAPI 240 segnalazioni si riducono a 46 voci. Le quattro categorie
tecniche hanno ora tutte un esito sulle coppie gemelle, pieno e senza falsi
positivi: nove casi su nove per l'Unrestricted Resource Consumption, tre su tre
per la Server-Side Request Forgery, sette su sette per la Security
Misconfiguration, tre su tre per l'Unsafe Consumption. La condizione di rischio
che emerge solo dalla congiunzione di segnali è quella dell'Unrestricted
Resource Consumption, dove un limite assente dal codice diventa una
segnalazione fondata solo se anche proxy e contratto non lo applicano. La
Su codice di terzi il quadro si divide. Sulla parte infrastrutturale Checkov,
con il catalogo delle policy, trova il 61\% delle misconfiguration documentate
di sadcloud; sulle API, nessuna delle voci delle quattro categorie tecniche
documentate da VAmPI e DVAPI viene trovata. La risposta a \RQ{1.1} è quindi
confermata sui casi di prova e, per la parte infrastrutturale, anche su codice
di terzi; per le categorie tecniche delle API non lo è ancora.

\paragraph{Risposta a \RQ{1.2}}
Nelle categorie ibride la parte statica riconosce la condizione necessaria,
cioè l'assenza del controllo, anche dove la verifica dinamica non
arriverebbe: lo mostrano i quattro veri positivi della BOLA su operazioni che il
contratto dichiara protette e il metodo del repository esterno assente dalla
specifica. Il residuo è di natura dinamica, e la campagna lo mostra nel modo più
diretto. Le due modalità di VAmPI hanno lo stesso codice: i livelli statici le
trattano allo stesso modo e solo il livello dinamico della BOLA le distingue.
Sul repository etichettato il livello dinamico conferma tre dei quattro metodi
che lo statico segnala e nessuno sulla risorsa protetta, mentre nella BOPLA la
segnalazione S04, vera sul contratto, è smentita dal sistema in esecuzione.
Nella BFLA il residuo è dinamico ma resta aperto: invocare la funzione è
facile, ma un \texttt{200 OK} non dice se l'invocazione fosse legittima, e le
segnalazioni BF-006 sugli endpoint di popolamento mostrano che a deciderlo è chi
conosce lo scopo della funzione, un'informazione che il criterio esclude di
fornire all'analisi. Nella Broken Authentication, infine, il residuo è dinamico
e condizionato: le prove distinguono le applicazioni gemelle, ma accertano un
difetto solo dove il meccanismo esiste e una sessione è disponibile.

\paragraph{Risposta a \RQ{1}}
Delle dieci categorie, quattro si comportano come tecniche (API4, API7, API8 e
API10) e si riconoscono sugli artefatti, come la IaC misconfiguration; quattro
sono ibride (API1, API2, API3 e API5) e richiedono in aggiunta un'evidenza di
natura dinamica, ottenuta sollecitando il sistema con identità diverse; API6 è
di logica di business, e API9 resta fuori dalla classificazione perché
presuppone un sistema non noto.

L'obiettivo della tesi è quindi raggiunto in parte. Lo è sul piano della
classificazione, che trova conferma dove i moduli sono stati provati: sulle
coppie gemelle ogni modulo distingue la versione vulnerabile da quella
corretta, e sull'insieme di prova la differenza fra le due modalità di VAmPI
emerge esattamente dove la classificazione la colloca, nel livello dinamico di
una categoria ibrida. Non lo è ancora sul piano della generalità: su crAPI i
livelli statici non segnalano nessuna delle sfide documentate, soprattutto per
il linguaggio e la forma del codice, e sull'insieme di prova la recall nel
perimetro è del 40\% su VAmPI e del 17\% su DVAPI, con una sola debolezza delle
API riconosciuta per quello che è, contro il 61\% sulla parte infrastrutturale
di sadcloud. La classificazione indica fin dove lo strumento può arrivare; il
divario fra quel limite e ciò che lo strumento raggiunge oggi su codice di terzi
è il risultato più importante della validazione, e il punto da cui partire per
estenderlo.

\section{Limiti e minacce alla validità}
\label{sec:ris-limiti}
\label{sec:ris-minacce}

\paragraph{Presenza, non correttezza}
L'analisi statica accerta la presenza di un controllo, non la sua correttezza:
è il residuo che la classificazione assegna alle categorie ibride. Nella BOLA un
filtro sul campo sbagliato viene accettato come legame, e una funzione chiamata
\texttt{authorize} viene accettata senza esaminarne il contenuto. Nella BFLA una
riga che cita il ruolo senza bloccare l'esecuzione, come un messaggio di log,
viene scambiata per una protezione, e nel contratto qualsiasi schema di
sicurezza basta a sopprimere BF-001.

\paragraph{Nomi, forme d'uso e perimetro}
Identità, controlli e funzioni riservate sono riconosciuti attraverso nomi e
forme d'uso diffusi nei framework più comuni: ciò che se ne discosta può
sfuggire, e un percorso dal nome allarmante ma legittimo può essere segnalato a
torto. Lo stesso vale per l'unità di analisi. Il livello statico della BOLA
esamina un file alla volta; BF-007 e l'Unrestricted Resource Consumption
seguono i dati solo dentro una funzione; la BFLA non analizza le viste definite
come classi e la BOPLA non analizza il codice Java. Per la Server-Side Request
Forgery restano fuori la forma cieca e i flussi che attraversano code
asincrone. Il caso di crAPI mostra quanto questi confini pesino su un
bersaglio realistico.

\paragraph{Configurazione}
Un controllo che non riesce a valutare un'espressione tace anziché fallire,
come \texttt{CKV\_AWS\_70}; il catalogo delle policy va aggiornato a ogni nuova
versione di Checkov; il manifest mostra ciò che dichiara, non ciò che è
attivo. Soprattutto, il flag di esposizione dice che una risorsa è
raggiungibile da chiunque, non se debba esserlo: un bucket che serve contenuti
pubblici viene segnalato come uno che espone dati riservati, e le segnalazioni
BF-006 ripropongono lo stesso limite sulle funzioni.

\paragraph{Parte dinamica e integrazione}
La parte dinamica richiede identità autenticate. Senza una sessione la Broken
Authentication si riduce alla componente statica, e la BOPLA, che non riceve
credenziali dal cruscotto né dalla scansione complessiva, lascia le proprie
segnalazioni da verificare. Il livello dinamico della BOLA richiede identità
reali, fornite da un bersaglio cooperante o da un traffico di esercizio i cui
token riportino il claim \texttt{sub}: DVAPI, che usa \texttt{userId}, ne
resta escluso. Sui bersagli esterni, privi di ripristino dello stato, il
livello dinamico prova inoltre i soli metodi che non modificano i dati, e le
violazioni che si manifestano solo con \texttt{PUT}, \texttt{PATCH} o
\texttt{DELETE} non vengono sollecitate. Il modello linguistico usato dalla
Broken Authentication introduce una fonte di non determinismo, che la campagna
ha reso visibile: sullo stesso repository etichettato il modulo ha prodotto una
segnalazione in un'esecuzione e nessuna nella successiva. L'integrazione è
ancora parziale: il rilevatore statico della BOLA non è registrato nei runner
della pipeline, e la Broken Authentication produce un proprio rapporto,
separato dal motore di correlazione.

\paragraph{Il prezzo del criterio}
Il criterio vieta di indicare all'analisi ciò che va protetto. Nella BOPLA una
proprietà è quindi protetta solo se l'applicazione la tratta come tale in
qualche punto, e un dato sensibile restituito a tutti senza alcun controllo di
ruolo non viene segnalato. È una conseguenza voluta: un elenco fisso di campi
sensibili avrebbe funzionato sulle applicazioni di prova e su nessun altro
sistema.

\paragraph{Minacce alla validità}
La minaccia principale alla validità interna è che casi di prova e regole hanno
lo stesso autore: i risultati pieni sulle coppie gemelle attestano la
conformità dell'implementazione alle intenzioni, non la sua efficacia su codice
reale. Le esecuzioni su codice esterno la attenuano solo in parte, perché
riguardano un repository per la BOLA e uno per la BFLA, e in entrambi il
riferimento di verità si è rivelato incompleto o ambiguo; il conteggio dei veri
e dei falsi positivi dipende quindi da che cosa si decide di considerare
vulnerabile. Il meccanismo che porta sulla stessa voce del motore di
correlazione conferme e smentite della BOLA è verificato solo dalle 68 prove
d'unità del livello dinamico, che simulano ZAP e le chiamate HTTP.

L'insieme di prova era pensato per eliminare la circolarità, ma questa prima
campagna la attenua soltanto. Lo strumento non è stato congelato prima
dell'esecuzione, come richiede il protocollo, e chi lo sviluppa ha visto i
risultati; durante la campagna sono stati corretti difetti dell'ambiente di
prova, non delle regole, fra cui il file di cattura del traffico condiviso fra
il bersaglio interno e quelli esterni, e le scansioni interessate sono state
ripetute. Le deviazioni sono annotate nel registro del protocollo, e la campagna
andrà ripetuta con lo strumento congelato.

Le metriche dipendono poi da scelte che conviene dichiarare. La regola di
corrispondenza guarda la categoria e il punto di esposizione, non la debolezza:
può contare come trovata una voce coperta da una segnalazione diversa sullo
stesso endpoint, come accade con due delle tre corrispondenze sulle API, e può
contare come mancata una debolezza rilevata sotto un'altra categoria, come
l'assenza di limiti di frequenza su VAmPI. Le decisioni manuali, cioè le
corrispondenze di sadcloud, le voci globali e il perimetro, sono state prese da
chi ha sviluppato lo strumento e sono registrate con il motivo, ma andrebbero
riviste da una persona esterna. La precisione, infine, è riportata solo come
limite inferiore, perché le segnalazioni che non corrispondono ad alcuna voce
non sono state classificate.

Quanto alla validità esterna, crAPI, VAmPI e DVAPI sono bersagli didattici,
costruiti perché le loro vulnerabilità siano riconoscibili, e crAPI cade in
larga parte fuori dal perimetro dei livelli statici: i risultati non si
estendono senza cautela ad applicazioni reali, né l'assenza di segnalazioni su
un bersaglio privo del difetto prova la sensibilità di un modulo. Le revisioni
fissate dei bersagli e la versione dichiarata di Checkov sostengono invece la
riproducibilità.
